% Section 2.5: the batch riders could add, r*/B = T0 X / M_kv, against the share
% of HBM the weights take, one curve per link bandwidth X. Reference GPU: 96 GB
% HBM at 4 TB/s, eta 0.75, 4 GB activation reserve; T0 = W/(eta BW),
% M_kv = 96 - W - 4. Computed 2026-10-20 (outline sec0-2, 2.5). 25.1 GB/s is
% the A100 PCIe4 copy rate we measure; 419 GB/s is Pie's measured GH200 C2C.
% 128 and 256 GB/s are plain bandwidths. TODO-verify: they correspond to CXL
% 3.x and 4.0 x16 links only once the CXL Consortium specification is cited.
% Target: Qwen2.5-72B fp8 on a GH200-96, w = 76/96 = 0.79, r*/B = 0.663.
\begin{figure}[t]
\centering
\begin{tikzpicture}
\begin{semilogyaxis}[
  width=0.86\columnwidth, height=5.0cm,
  xlabel={share of HBM taken by the weights},
  ylabel={riders per resident, $r_0/B$},
  xmin=0.28, xmax=0.92, ymin=0.003, ymax=4,
  xtick={0.3,0.4,0.5,0.6,0.7,0.8,0.9},
  xticklabels={30\%,40\%,50\%,60\%,70\%,80\%,90\%},
  ytick={0.01,0.1,1}, yticklabels={0.01,0.1,1},
  tick label style={font=\sffamily\scriptsize},
  label style={font=\sffamily\scriptsize},
  axis line style={black!60}, tick style={black!60},
  ymajorgrids, grid style={black!10},
  every axis plot/.append style={line width=0.9pt},
  clip=false,
]
\pgfplotstableread{
w     p4     p5     x128   x256   c2c
0.30  0.004  0.010  0.019  0.039  0.064
0.40  0.006  0.015  0.031  0.061  0.100
0.50  0.009  0.023  0.047  0.093  0.152
0.60  0.014  0.036  0.071  0.143  0.234
0.70  0.023  0.058  0.116  0.231  0.378
0.75  0.030  0.077  0.154  0.307  0.503
0.79  0.039  0.100  0.200  0.400  0.655
0.83  0.054  0.138  0.276  0.552  0.903
0.87  0.082  0.210  0.420  0.840  1.376
0.90  0.129  0.329  0.658  1.317  2.155
}\potential
\addplot[black!45]                 table[x=w, y=p4]   {\potential};
\addplot[black!70, densely dashed] table[x=w, y=p5]   {\potential};
\addplot[h2dC!60]                  table[x=w, y=x128] {\potential};
\addplot[h2dC, densely dashed]     table[x=w, y=x256] {\potential};
\addplot[hbmD, line width=1.2pt]   table[x=w, y=c2c]  {\potential};
\node[font=\sffamily\scriptsize, text=black!55, anchor=west] at (axis cs:0.905,0.129) {25\,GB/s};
\node[font=\sffamily\scriptsize, text=black!75, anchor=west] at (axis cs:0.905,0.329) {64};
\node[font=\sffamily\scriptsize, text=h2dC!70, anchor=west]  at (axis cs:0.905,0.62) {128};
\node[font=\sffamily\scriptsize, text=h2dC, anchor=west]     at (axis cs:0.905,1.2) {256};
\node[font=\sffamily\scriptsize, text=hbmD, anchor=west]     at (axis cs:0.905,2.45) {419};
\addplot[only marks, mark=*, mark size=2.2pt, draw=black, fill=ddrD]
  coordinates {(0.7917,0.663)};
\node[font=\sffamily\scriptsize, anchor=east, align=right] (tg)
  at (axis cs:0.70,1.7) {GH200, Qwen2.5-72B fp8};
\draw[black!60, line width=0.5pt] (tg.east) -- (axis cs:0.785,0.72);
\end{semilogyaxis}
\end{tikzpicture}
\caption{\textbf{What the link could add.} Requests the link can carry per step
relative to those HBM holds, $r_0/B = T_0\Xlink/\Mkv$, for a GPU with
96\,GB of HBM at 4\,TB/s, one curve per link bandwidth $\Xlink$ (GB/s; 25 and
419 are measured rates). An upper bound, before the costs $F$ and $c$ of
\S\ref{sec:model}.}
\label{fig:potential}
\end{figure}
